Penelitian ini dilaksanakan dengan menjunjung tinggi prinsip etika penelitian dan perlindungan terhadap komunitas Tuli. 
BISINDO dipandang sebagai bahasa alami dan bagian dari identitas budaya komunitas Tuli, sehingga penggunaannya dalam penelitian harus dilakukan secara bertanggung jawab.

Prinsip etika yang diterapkan dalam penelitian ini meliputi:
\begin{enumerate}[noitemsep]
    \item Tidak melakukan perekaman data tanpa persetujuan yang jelas dan tertulis dari pihak yang terlibat.
    \item Menghormati hak dan martabat penutur Tuli, serta tidak memposisikan mereka sebagai objek uji coba.
    \item Membatasi penggunaan data hanya untuk kepentingan akademik dan penelitian.
    \item Tidak mengklaim bahwa sistem yang dikembangkan bertujuan menggantikan peran guru, penerjemah, atau interaksi manusia dalam pembelajaran BISINDO.
\end{enumerate}

Pendekatan etis ini diterapkan untuk memastikan bahwa penelitian tidak hanya valid secara teknis, tetapi juga bertanggung jawab secara sosial dan budaya.
